\documentclass[10pt,a4paper]{amsart}
\usepackage[margin=19mm]{geometry}
\usepackage{amsmath,amssymb,mathtools}
\usepackage[scr=rsfs,cal=euler]{mathalfa}
\usepackage{xfrac}
\usepackage[unicode,hidelinks,bookmarks=false]{hyperref}
\newcommand{\ie}{i.e.\ }
\newcommand{\cf}{cf.\ }
\newcommand{\resp}{respectively\ }
\newcommand{\Subsection}{\subsection}
\providecommand{\nobreakdash}{\nobreak\hbox{-}}
\setlength{\emergencystretch}{2em}
\multlinegap0pt
\usepackage{bm}
\usepackage{bbm}
\usepackage{dsfont}
\usepackage{xspace}
\usepackage{cancel}
\usepackage{mathtools}
\usepackage{amsthm}
\makeatletter
\@ifundefined{Theorem}{\newtheorem{Theorem}{Theorem}}{}
\@ifundefined{Lemma}{\newtheorem{Lemma}[Theorem]{Lemma}}{}
\makeatother
\def\H{\mathcal H}
\def\R{\mathbb R}
\def\d{\mathrm{d}}
\def\e{\mathop{\mathrm{e}}\nolimits}
\title[Levinson's theorem for dissipative systems, or how to count ]{Levinson's theorem for dissipative systems, or how to count the asymptotically disappearing states}
\author{A. Alexander, J. Faupin, S. Richard}
\date{Source: arXiv:2509.12799v1 (CC BY 4.0)}
\begin{document}
\maketitle

\noindent\emph{Source and licence.} Levinson's theorem for dissipative systems, or how to count the asymptotically disappearing states. Version arXiv:2509.12799v1 (CC BY 4.0), distributed under the Creative Commons Attribution 4.0 International licence, as recorded in the arXiv licence metadata for this exact version and shown on the abstract page https://arxiv.org/abs/2509.12799v1. Full rights notice: licenses/arxiv-2509.12799-CC-BY-4.0.txt.\par\smallskip

\noindent\textit{Editorial scope.} Counted unit: the Lemma whose source label is \texttt{lem\_statio} in the article (source0.tex lines 396--409, source label \texttt{lem\_statio}), delivered together with the article's own complete original proof of it (source0.tex lines 411--448, one \texttt{proof} environment of 38 lines). No mathematics is written, repaired or re-derived: statement and proof are the article's own text line for line. Required context retained uncounted: none. Every definition, display and label of those lines is retained unchanged. Omitted from this package and not redistributed: 1--395, 410--410, 449--1694. Nothing inside a retained excerpt is omitted or altered: every retained body is delivered exactly as the article's lines are written, so no omission receipt of any kind accompanies this package. Citations: the retained excerpts carry no citation command at all, so no bibliography entry of the article is transcribed and no reference is redistributed. Reference adaptation: none was needed and none was made; every reference inside the delivered unit resolves inside it, so no unresolved reference or citation is produced. Printed slips kept literal: no printed word, symbol, hypothesis or number of the counted statement or of its proof was altered. Layout and class: the article's own class is \texttt{article} with packages including times, amssymb, amsmath, amsthm, comment, mathrsfs, epsfig, color, enumitem; the wrapper is the lane's pinned \texttt{amsart} editorial wrapper at 10pt on A4, which restates the theorem-like environments the retained text needs and adds the article's own macro definitions that the retained text needs (\texttt{\textbackslash H}, \texttt{\textbackslash R}, \texttt{\textbackslash d}, \texttt{\textbackslash e}), with extra wrapper packages (bm, bbm, dsfont, xspace, cancel, mathtools). The delivered numbering is the unit's own. Assets: no retained span contains a figure environment, a graphics-inclusion command, a PDF-page inclusion, a table or a TikZ picture; no image file is copied into this package, the package declares no asset dependency and no image file is redistributed with it. Licence: version arXiv:2509.12799v1 of this article is distributed under the Creative Commons Attribution 4.0 International licence, as recorded in the arXiv licence metadata for this exact version and shown on the abstract page https://arxiv.org/abs/2509.12799v1. A full rights notice is delivered at licenses/arxiv-2509.12799-CC-BY-4.0.txt. Characters: every retained excerpt is pure ASCII, so the delivered engine, XeLaTeX with its default Latin Modern text font, reports no missing character.
\begin{Lemma}\label{lem_statio}
If the wave operators $W_\pm$ exist, then for any $f, g\in \H$ the following equalities hold:
$$
\langle W_+ f,g\rangle 
=  \lim_{\varepsilon\searrow 0}\frac{\varepsilon}{\pi} \int_{\R}\d \lambda 
\langle R_0(\lambda+i \varepsilon)f, R(\lambda + i \varepsilon) g\rangle
$$
and 
$$
\langle W_- f,g\rangle
=  \lim_{\varepsilon\searrow 0}\frac{\varepsilon}{\pi} \int_{\R}\d \lambda 
\langle R_0(\lambda- i \varepsilon)f, R^*(\lambda - i \varepsilon) g\rangle.
$$
\end{Lemma}

\begin{proof}
We give the proof for $W_+$, the proof for $W_-$ is similar. 
Since $\{ \e^{-itH} \}_{ t \ge 0 }$ is a contraction semigroup and $\{ \e^{-itH_0} \}_{ t \in\R }$ is a unitary group, we have 
\begin{equation*}
R(\lambda+i\varepsilon)f=i\int_0^\infty \d t\, \e^{-\varepsilon t}\e^{-it(H-\lambda)}f, \qquad R_0(\lambda+i\varepsilon)f=i\int_0^\infty \d t\, \e^{-\varepsilon t}\e^{-it(H_0-\lambda)}f,
\end{equation*}
for all $\varepsilon>0$ and $f\in\H$ with $\|f\|=1$, and 
\begin{equation*}
\frac{\varepsilon}{\pi} \int_{\R}\d \lambda \|R(\lambda+i\varepsilon)f\|^2\le1, \qquad \frac{\varepsilon}{\pi} \int_{\R}\d \lambda \|R_0(\lambda+i\varepsilon)f\|^2=1.
\end{equation*}
Let $\varepsilon>0$ and $f,g\in\H$. Using Lebesgue's dominated convergence Theorem, we write
\begin{align}\label{eq:int-formula}
&\frac{\varepsilon}{\pi} \int_{\R}\d \lambda \langle R_0(\lambda+i \varepsilon)f, R(\lambda + i \varepsilon) g\rangle =\frac{\varepsilon}{\pi}\lim_{\delta\searrow 0}  \int_{\R}\d \lambda \e^{-\delta\lambda^2} \langle R_0(\lambda+i \varepsilon)f, R(\lambda + i \varepsilon) g\rangle , 
\end{align}
and then compute, for all $\delta>0$,
\begin{align*}
&\int_{\R}\d \lambda \e^{-\delta\lambda^2} \langle R_0(\lambda+i \varepsilon)f, R(\lambda + i \varepsilon) g\rangle \\
&= \int_{\R}\d \lambda \int_0^\infty \d t \int_0^\infty \d t' \e^{-\delta\lambda^2} \e^{-\varepsilon (t+t')} \e^{i\lambda(t'-t)} \langle \e^{-itH_0}f, \e^{-it'H} g\rangle \\
&= \int_{\R}\d \lambda \int_0^\infty \d t \int_{-t}^\infty \d s \,  \e^{-\delta\lambda^2} \e^{-2\varepsilon t} \e^{-\varepsilon s} \e^{i\lambda s} \langle \e^{-itH_0}f, \e^{-i(t+s)H} g\rangle ,
\end{align*}
where we used the change of variables $t'=t+s$ in the last equality. Using Fubini's Theorem and computing the integral in $\lambda$, we obtain
\begin{align*}
&\int_{\R}\d \lambda \e^{-\delta\lambda^2} \langle R_0(\lambda+i \varepsilon)f, R(\lambda + i \varepsilon) g\rangle \\
&= (2\pi)^\frac12 (2\delta)^{-\frac12} \int_0^\infty \d t \int_{-t}^\infty \d s \,  \e^{-s^2/(4\delta)} \e^{-2\varepsilon t} \e^{-\varepsilon s}  \langle \e^{-itH_0}f, \e^{-i(t+s)H} g\rangle \\
&= (2\pi)^{\frac12} 2^\frac12 \int_0^\infty \d t \int_{-(4\delta)^{-1/2}t}^\infty \d \tau \,  \e^{-\tau^2} \e^{-2\varepsilon t} \e^{-(4\delta)^{1/2}\varepsilon \tau}  \langle \e^{-itH_0}f, \e^{-itH} \e^{-i(4\delta)^{1/2}\tau H} g\rangle ,
\end{align*}
where we used the change of variables $\tau=(4\delta)^{-1/2}s$. Applying again Lebesgue's dominated convergence Theorem and then computing the integral in $\tau$, we obtain
\begin{align*}
\lim_{\delta\searrow 0}\int_{\R}\d \lambda \e^{-\delta\lambda^2} \langle R_0(\lambda+i \varepsilon)f, R(\lambda + i \varepsilon) g\rangle &= 2\pi \int_0^\infty \d t \e^{-2\varepsilon t}  \langle \e^{-itH_0}f, \e^{-itH} g\rangle.
\end{align*}
Inserting this into \eqref{eq:int-formula} yields
\begin{align*}
\frac{\varepsilon}{\pi} \int_{\R}\d \lambda \langle R_0(\lambda+i \varepsilon)f, R(\lambda + i \varepsilon) g\rangle &= 2\varepsilon \int_0^\infty \d t \e^{-2\varepsilon t}  \langle \e^{-itH_0}f, \e^{-itH} g\rangle\\
&= 2 \int_0^\infty \d s \e^{-2s}  \langle \e^{-i\varepsilon^{-1}sH_0}f, \e^{-i\varepsilon^{-1}sH} g\rangle \\
&= 2 \int_0^\infty \d s \e^{-2s}  \langle \e^{i\varepsilon^{-1}sH^*} \e^{-i\varepsilon^{-1}sH_0}f,  g\rangle ,
\end{align*}
by the change of variables $s=\varepsilon t$. A last application of Lebesgue's dominated convergence Theorem gives the result.
\end{proof}

\end{document}
